\section{Introduction}\label{sec:intro}

An $m$-variate polynomial over the finite field~$\F_q$ is a function $g:\F_q^m \rightarrow \F_q$ of the form
\begin{equation*}
g(x_1, \ldots, x_m) = \sum_{i_1, \ldots, i_m} c_{i_1, \ldots, i_m} \cdot x_1^{i_1} \cdots x_{m}^{i_m},
\end{equation*}
where each coefficient~$c_{i_1, \ldots, i_m}$ is an element of~$\F_q$.
We say that~$g$ has \emph{total degree $d$}
(or \emph{degree $d$}, for short) if $i_1 + \cdots + i_m \leq d$ for each nonzero coefficient $c_{i_1, \ldots, i_m}$,
and \emph{individual degree $d$} if $i_1, \ldots, i_m \leq d$ for each nonzero coefficient $c_{i_1, \ldots, i_m}$.
Low-degree polynomials have a variety of properties which make them useful in theoretical computer science,
chief among which is their \emph{distance}:
by the Schwartz-Zippel lemma, two nonequal degree $d$ polynomials~$g$ and~$h$ agree on at most a $d/q$ fraction of the points in~$\F_q^m$.



\emph{Low (individual) degree testing} refers to the task 
of verifying that an unknown function $g:\F_q^m \rightarrow \F_q$
is representable as a polynomial of (individual) degree~$d$
by querying~$g$ on a small number of points $u \in \F_q^m$.
There is a pair of canonical tests for doing so
known as the \emph{surface-versus-point low-degree test}
and the \emph{low individual degree test}.
It is common to frame these tests
as games between a referee and two provers.
In this setting, the surface-versus-point low degree test, parameterized by an integer $k \geq 1$, is performed by the verifier as follows.
\begin{enumerate}
\item Select~$\bu \sim \F_q^m$ uniformly at random. Give it to Prover~$\mathrm{A}$. They respond with a value~$\ba \in \F_q$.
\item Select a uniformly random $k$-dimensional affine surface~$\bs$ in~$\F_q^m$ containing~$\bu$.
	Give it to Prover~$\mathrm{B}$. They respond with a degree-$d$ $k$-variate polynomial $\boldf:\bs\rightarrow \F_q$.
\item Accept if $\boldf(\bu) = \ba$.
\end{enumerate}
This test is motivated by the following ``local characterization" of low-degree polynomials:
a polynomial $g:\F_q^m \rightarrow \F_q$ is degree-$d$
if and only if $g|_s$ is degree-$d$ for all $k$-dimensional surfaces~$s$.
Hence, if~$g$ is degree-$d$,
then the provers can win with probability~$1$ by always replying with~$\ba = g(\bu)$ and $\boldf = g|_{\bs}$.
\emph{Soundness} of the low-degree test refers to the converse statement,
namely that players who succeed with high probability
must be responding based on a low-degree polynomial.
This is formalized as follows.

\begin{theorem}[Raz-Safra~\cite{RS97}]\label{thm:raz-safra}
Suppose Provers~$\mathrm{A}$ and~$\mathrm{B}$ pass the $k = 2$ surface-versus-point low-degree test with probability~$1-\eps$.
Then there exists a degree-$d$ polynomial $g:\F_q^m \rightarrow \F_q$ such that
\begin{equation*}
\Pr_{\bu \sim \F_q^m}[g(\bu) = \ba] \geq 1 - \eps - \poly(m) \cdot \poly(d/q).
\end{equation*}
\end{theorem}

A similar ``local characterization" of low individual degree polynomials states that a polynomial~$g$ has individual degree $d$
if and only if $g|_{\ell}$ is a univariate degree-$d$ polynomial for all axis-parallel lines~$\ell$.
An axis-parallel line is a line of the form $\ell = \{u + a \cdot e_i \mid i \in \F_q\}$, for $u \in \F_q^m$ and $i \in \{1, \ldots, m\}$.
Motivated by this, the low individual degree test follows the same outline as the surface-versus-point low degree test
except with the second step substituted with the following.
\begin{enumerate}
\setcounter{enumi}{1}
\item Select a uniformly random axis-parallel line $\bell$ in~$\F_q^m$ containing~$\bu$.
	Give it to Prover~$\mathrm{B}$. They respond with a degree-$d$ univariate polynomial $\boldf:\bell\rightarrow \F_q$.
\end{enumerate}
When $d = 1$, the low individual degree test
is called the \emph{multilinearity test} because a polynomial with individual degree $d = 1$ is a multilinear polynomial.
The multilinearity  and low individual degree tests were
first introduced and proven sound by Babai, Fortnow, and Lund in~\cite{BFL91}.
The analysis of its soundness was then improved by~\cite{AS98} and then further sharpened by~\cite{FHS94}.
The best bound follows from the work of Polishchuk and Spielman~\cite{PS94};
their work considers only the bivariate $m = 2$ case,
but extending it to the multivariate case yields the following result.
\begin{theorem}[Polishchuk-Spielman~\cite{PS94}]\label{thm:classical-test-soundness}
Suppose Provers~$\mathrm{A}$ and~$\mathrm{B}$ pass the low individual degree test with probability~$1-\eps$.
Then there exists a polynomial $g:\F_q^m \rightarrow \F_q$ with individual degree~$d$ such that
\begin{equation*}
\Pr_{\bu \sim \F_q^m}[g(\bu) = \ba] \geq 1 -  \poly(m) \cdot (\poly(\eps) + \poly(d/q)).
\end{equation*}
\end{theorem}
We note that the soundness error the low individual test gives is actually worse than the low degree test,
because the low individual degree function~$g$ is only $\poly(m) \cdot (\poly(\eps) + \poly(d/q))$ close to Player~$\mathrm{A}$'s strategy,
rather than $\eps + \poly(m) \cdot \poly(d/q)$.
We will discuss this weakness of the low individual degree test below.

The multilinearity test, low individual degree test, and low-degree test
form a sequence
in which each test generally enables more applications than the previous one.
The multilinearity test can be used to show that $\MIP = \NEXP$ using a polynomial number of rounds~\cite{BFL91},
the low individual degree test can reduce the number of rounds to~$1$,
and the low degree test can be used to ``scale this result down''
and prove the PCP theorem, i.e.\ $\NP = \MIP[O(\log(n)), O(1)]$~\cite{AS98,ALM+98}.

\subsection{Quantum soundness of the low degree tests}

The work of Ito and Vidick~\cite{IV12}
initiated a program of studying these tests
in the case when the players are quantum,
as a means of proving bounds on the complexity class~$\MIP^*$.
Because the provers are quantum, 
they are allowed to share an entangled state,
a resource which could potentially allow them to ``cheat" the test
and win without using a low-degree polynomial.
The goal of this program is to show that this is not possible.
In other words, the goal is to show that these tests are \emph{quantum sound},
which means that provers who succeed with high success probability
must answer their questions according to a low (individual) degree polynomial,
even if they are allowed to share quantum entanglement.\footnote{We note that \emph{quantum soundness of the low-degree tests}, which is the focus of this work,
is distinct from \emph{soundness of the quantum low-degree test}.
The ``quantum low-degree test'' is a particular test introduced by Natarajan and Vidick in~\cite{NV18a}
which gets its name from the prominent role that the low-degree test plays as a subroutine,
and its ``soundness'' is simply the result that they prove about it.
}

Correctly formalizing the notion of quantum soundness is a subtle task,
as quantum provers can in fact ace these tests
using a broader class of strategies than their classical counterparts.
For example, the two provers can use their quantum state~$\ket{\psi}$ to simulate shared randomness,
which they can use to sample a random low-degree polynomial~$\bg$  to answer their questions with;
what makes this still acceptable is that~$\bg$ depends only on their shared randomness and not their questions.
The correct formalization of quantum soundness was identified by Ito and Vidick~\cite{IV12},
which states the following:
suppose Provers~$\mathrm{A}$ and~$\mathrm{B}$ pass the low-degree test with probability close to~$1$.
For each point question~$u \in \F_q^m$,
let $A^{u} = \{A^{u}_a\}$ be the measurement that Prover~$\mathrm{A}$ applies to
their share of~$\ket{\psi}$ to produce the answer $\ba \in \F_q$.
Then the test being quantum sound
means that there should be a measurement $G = \{G_g\}$, independent of~$u \in \F_q^m$,
which outputs degree-$d$ polynomials~$g$
and ``acts like~$A$''.
In other words, rather than measuring~$A^{u}$ to produce the outcome~$\ba \in \F_q$,
Prover~$\mathrm{A}$ could have simply measured~$G$, received the polynomial~$\bg$,
and outputted its evaluation at~$\bu$, i.e.\ the value $\bg(\bu)$.
We will measure the similarity between~$A$ and~$G$
by considering the experiment where Prover~$\mathrm{A}$ measures with~$A^{\bu}$ to produce~$\ba$,
Prover~$\mathrm{B}$ measures with~$G$ to produce~$\bg$, and we check if $\bg(\bu) = \ba$.
This entails studying the quantity
\begin{equation*}
\E_{\bu \sim \F_q^m} \sum_{a \in \F_q} \sum_{g:g(\bu) = a}\bra{\psi} A^{\bu}_a \ot G_g \ket{\psi},
\end{equation*}
which we aim to show is as close to~$1$ as possible.
In this way, the provers' quantum advantage is limited to their ability to select a low-degree polynomial~$g$.

One additional quirk of the quantum setting
is that it has been historically useful to consider variants of these tests which feature more than two provers.
This allows one to use monogamy of entanglement to reduce the power that entanglement gives to the provers,
making it easier to show that a given test is quantum sound.
Low degree test results with fewer provers are more difficult to show and have more applications.

The program of showing that these tests are quantum sound was carried out for the $3$-prover multilinearity test by Ito and Vidick~\cite{IV12}
and for the $3$-prover low-degree test by Vidick~\cite{Vid16},
which was later improved to $2$-provers by Natarajan and Vidick~\cite{NV18b}.
This latter result led to a sequence of works
which culminated in the proof that $\mathsf{MIP}^* = \mathsf{RE}$
and the refutation of the Connes embedding conjecture in~\cite{JNV+20}.
We summarize this line of research in \Cref{fig:research}.



{
\floatstyle{boxed} 
\restylefloat{figure}
\begin{figure}
\begin{tabular}{l c c c}
& Test shown & Complexity-theoretic  & Number \\
&  quantum-sound &  consequence & of provers  \\
			& &&\\
1. \cite{IV12}: & multilinearity test & $\NEXP \subseteq \mathsf{MIP}^*$ & 3\\[1ex]
2. \cite{Vid16}: & low-degree test & $\NP \subseteq \mathsf{MIP}^*[O(\log(n)), O(1)]$  &3\\[1ex]
3. \cite{NV18b}: & low-degree test & $\NP \subseteq \mathsf{MIP}^*[O(\log(n)), O(1)]$ & 2\\
			& &&\\
& Consequences of \cite{NV18b}: & &\\
			& & &\\
&\multicolumn{1}{l}{\quad\qquad(a) \cite{NV18a}:}   & $\QMA \subseteq \mathsf{MIP}^*[O(\log(n)), O(1)]$ &7\\
&    & (under randomized reductions) &\\[1ex]
&\multicolumn{1}{l}{\quad\qquad (b) \cite{NW19}:}   & $\NEEXP \subseteq \mathsf{MIP}^*$ &2\\[1ex]
&\multicolumn{1}{l}{\quad\qquad (c) \cite{JNV+20}:}   & $\MIP^* = \mathsf{RE}$&2
\end{tabular}
	\caption{Prior work on quantum-sound low degree tests and their complexity-theoretic consequences.
			The first three works showed a quantum-sound test and an $\MIP^*$ bound,
			both involving the same number of provers indicated in the final column.
			The last three works use the low-degree test from~\cite{NV18b}
			to show the indicated $\MIP^*$ bound.}
\label{fig:research}
\end{figure}
}


Subsequent to initial posting of~\cite{JNV+20} on the arXiv,
an error was discovered in the analysis of the quantum-sound low degree test
contained in~\cite{Vid16} which was propagated to~\cite{NV18b}. The error affects the proof in a manner that appears difficult to fix. As such, we currently do not know if the low-degree test is quantum-sound for any number of provers.
The invalidation of this analysis affects every result in \Cref{fig:research} except for~\cite{IV12}.


The purpose of this work is to provide a different soundness analysis, for a variant of the low-degree test, that can nevertheless be used as a replacement for it in most subsequent works. 
We do so by revisiting the three-player quantum-sound multilinearity test of~\cite{IV12}
and improving this result in two ways.
First, we generalize it to hold for the degree-$d$ low individual degree test,
of which the multilinearity test is the $d=1$ special case.
Second, using techniques introduced in~\cite{Vid16,NV18b},
we reduce the number of provers from~$3$ to~$2$.
Our main result is as follows.

\begin{theorem}[Main theorem, informal]\label{thm:main-informal}
Suppose Provers~$\mathrm{A}$ and~$\mathrm{B}$
pass the two-prover. degree-$d$ low individual degree test with probability $1-\eps$.
Let $A = \{A^u_a\}$ be the measurement the provers perform when they are given the point $u \in \F_q^m$
to produce a value $a \in \F_q$.
Then there exists a projective measurement $G = \{G_g\}$
whose outcomes~$g$ are polynomials of individual degree~$d$
such that
\begin{equation*}
\E_{\bu \sim \F_q^m}\sum_{a \in \F_q} \sum_{g:g(\bu) = a} \bra{\psi} A^{\bu}_{a} \ot G_g \ket{\psi}
\geq 1 - \poly(m) \cdot (\poly(\eps) + \poly(d/q)).
\end{equation*}
In other words, if Prover~$\mathrm{A}$ measures according to~$A^{\bu}$ to produce $\ba$ and Prover~$\mathrm{B}$ measures according to~$G$ to produce~$\bg$,
then $\bg(\bu) = \ba$ except with probability $\poly(m) \cdot (\poly(\eps) + \poly(d/q))$.
\end{theorem}
\noindent
Thus, we are able to extend \Cref{thm:classical-test-soundness}
to the case of two quantum provers (with some minor caveats; see \Cref{thm:main-formal} below for the formal statement of \Cref{thm:main-informal}).

Although \Cref{thm:main-informal} establishes quantum soundness of the low-individual degree test
and not of the low-degree test, it is still sufficient to recover the result $\NEEXP \subseteq \mathsf{MIP}^*$
from~\cite{NW19}
and the result $\mathsf{MIP}^* = \mathsf{RE}$ from~\cite{JNV+20}.
In addition,
we can use it to recover the self-test for an exponential number of EPR pairs from~\cite{NV18a}.
Edited drafts of these work to account for this change are forthcoming.
It remains open whether the complexity-theoretic consequences
of~\cite{Vid16,NV18b,NV18a} to the ``scaled down'' setting still hold.

\subsection{Total degree versus individual degree}

We now contrast the low degree test with the individual degree test
and explain why we are only able to prove the latter quantum sound.
We begin by explaining why the low individual degree test,
unlike the low degree test,
requires a $\poly(m) \cdot \poly(\eps)$ dependence in the soundness error.

\begin{example}\label{ex:bad-individual-degree-example}
Consider the degree-$(d+1)$ polynomial $h(x_1, \ldots, x_m) = x_1^{d+1}$,
and suppose that Players~$\mathrm{A}$ and~$\mathrm{B}$ play according to the following classical strategy.
\begin{itemize}
\item[$\circ$] (Player~$\mathrm{A}$): given $\bu \in \F_q^m$, return the value $\ba = h(\bu)$.
\item[$\circ$] (Player~$\mathrm{B}$): given the axis parallel line~$\bell = \{\bu + x \cdot e_{\bi} \mid x \in \F_q\}$, act as follows.
				If $\bi > 1$, then $h$ is a constant function along~$\bell$, and so return $\boldf = h|_{\bell}$.
				Otherwise, if $\bi = 1$, then $h$ is a degree-$(d+1)$ polynomial along~$\bell$.
				As the verifier expects a degree-$d$ polynomial,
				simply give up and return $\boldf \equiv 0$.
\end{itemize}
Using this strategy, $\boldf(\bu) = h|_{\bell}(\bu) = h(\bu) = \ba$ whenever $\bi \neq 1$, which occurs with probability $1- \frac{1}{m}$.
Hence, the players pass the degree-$d$ low individual degree test with probability at least $1 -\eps$ for $\eps = \frac{1}{m}$.
However, Player~$\mathrm{A}$ is responding to their questions using the polynomial~$h$
which is degree-$(d+1)$ but not degree-$d$,
and so by the aforementioned Schwartz-Zippel lemma,
any degree-$d$ polynomial~$g$ will agree with~$h$ on at most a $\frac{d+1}{q}$ fraction of all inputs.
This means that the agreement between Player~$\mathrm{A}$'s strategy and any degree-$d$ polynomial~$g$ is at most
\begin{equation*}
\Pr_{\bu \sim \F_q^m}[g(\bu) = \ba]
\leq 
1 - m \cdot \eps + \frac{d+1}{q}.
\end{equation*}
\end{example}

\Cref{ex:bad-individual-degree-example} shows that the dependence on~$m$ and~$\eps$ in \Cref{thm:main-informal} is tight up to polynomial factors.
This reveals a weakness with the low individual degree test:
one can only conclude that the players are using a low individual degree strategy
when their failure probability $\eps$ is tiny---on the order of $\frac{1}{m}$ or smaller.
The low (total) degree test is alluring because it has the potential to avoid this dependence on~$m$.

Soundness of the low total and individual degree tests
is typically proven by induction on~$m$.
For the low individual degree test,
each step of the induction incurs an error of $\poly(m) \cdot (\poly(\eps) + \poly(d/q))$.
Summing over all~$m$ steps, this gives a total error of $\poly(m) \cdot (\poly(\eps) + \poly(d/q))$,
exactly as in \Cref{thm:main-informal}.
For the low degree test,
on the other hand, each step of the induction only incurs an error of $\poly(\eps) + \poly(d/q)$.
However, if summed over all~$m$ steps,
this still gives a total error of $m \cdot (\poly(\eps) + \poly(d/q))$,
which is too large.

To account for this,
the soundness proofs in \cite{Vid16, NV18b}
introduce a technique at the end of each induction step
called \emph{Consolidation}
in which the growing error is ``reset''
down to an error $\poly(\eps) + \poly(d/q)$
which remains fixed across all levels of the induction.
This allows them to conclude with an error that was independent of the dimension~$m$.
Consolidation works as follows:
if the error of the projective measurement $G = \{G_g\}$
ever grows past this fixed error, the Consolidation step argues that on some portion of the Hilbert space,
the measurement~$G$ must be performing much worse than expected;
it then corrects this by inductively calling the low-degree test soundness to produce a better measurement on this portion of the Hilbert space.
Ultimately, however, this creates a cascade of Consolidation steps calling each other with increasing error,
which at some point grows so large that the low-degree soundness can no longer be applied.
This is the source of the bug.



Fortunately, this work shows that the techniques of these prior works, 
aside from Consolidation,
are still sound.
So we can show soundness bounds which grow as a function of~$m$,
even if we cannot yet show bounds independent of~$m$.
In the end, this ``weakness" of the low individual degree test
is precisely what allows us to show that it, and not the low degree test, is quantum sound.

That said, we do believe, although we have not rigorously verified,
that our techniques are capable of proving a result like \Cref{thm:main-informal} for the low degree test,
i.e.\ a soundness bound of $\poly(m) \cdot (\poly(\eps) + \poly(d/q))$ rather than $\eps + \poly(m) \cdot \poly(d/q)$.
This ``weak'' quantum soundness does not rule out the possibility that quantum provers can significantly outperform their classical counterparts.
It is, however, still sufficient for the same applications as the low individual degree test,
and it hints towards the possibility of a full quantum soundness of the low degree test.

\subsection{Conclusion and open problems}

In recent years,
the classical low-degree test has played a critical role in the study of nonlocal games
and the complexity class~$\MIP^*$.
In addition to correcting previous proofs of soundness,
we hope that this new exposition
will invite new researchers to engage with this beautiful area.
We conclude with a short list of open problems for future work.

\begin{enumerate}
\item Is the classical low-degree test quantum sound? Can the Consolidation step be fixed?
\item Even if the classical low-degree test is shown quantum sound,
	there are still interesting questions to be answered about the low individual degree test.
	For example, can the diagonal lines test be removed? (See \Cref{sec:the-test} for a description of this subtest.)
	Doing so would likely simplify the proof of $\MIP^* = \mathsf{RE}$~\cite{JNV+20},
	as one could replace the complicated ``conditional linear functions'' used in the proof
	with a simpler subclass known as ``coordinate deletion functions''.
	However, as discussed at the end of \Cref{sec:technical-overview},
	we know of an example that requires the diagonal lines test
	for the low individual degree test with parameters $m = 2$, $d = 2$, and $q = 4$.
	Can we find similar examples for larger~$q$?
\item Classically,  the low degree test
	generalized in various directions,
	including tests for affine invariant properties~\cite{KS08,Sud11}
	and tests for tensor product codes (see, for example, \cite{CMS17}).
	Does quantum soundness hold for these tests as well?
\item This work continues the trend of showing that well-studied classical property testers
	are also quantum sound.
	Prior to this, the work of~\cite{IV12} (see also \cite[Chapter 2]{Vid11}) showed that the linearity tester of Blum, Luby, and Rubinfield~\cite{BLR93}
	is also quantum sound.
	The proofs of these results have so far been case-by-case adaptations of the classical proofs to the quantum setting;
	in the case of this paper, the proof is highly nontrivial
	and involves many ad hoc calculations which fortunately go in our favor.
	Is there a more conceptual reason why quantum soundness holds for these testers?
	Perhaps a reduction from the quantum case to the classical case?
\end{enumerate}


% \section{Introduction old}

% Self-testing is a big deal in quantum information theory and lies at
% the heart of the recent work on MIP*. But actually the term
% ``self-testing'' originates in the \emph{classical} computer science
% literature.  Blum, Luby, and Rubinfeld define a program to compute a
% function as ``self-testable'' if it can be self-corrected

% One of the main technical tools BLR developed was a test for linear
% functions. Given query access to a function $f: \{0,1\}^n \to \{0,1\}$, the
% BLR test certifies that $f$ is close to some \emph{linear} function
% $g(x ) = a \cdot x$ by making queries to $f$ at three randomly chosen
% points. To show this, BLR demonstrate how to ``self-correct'' any
% function $f$ that passes the test with high probability to obtain a
% linear $g$ that agrees with $f$ on most points. Moreover, BLR showed a
% stronger property: an agent given only query access to $f$ can with
% high probability
% evaluate $g$ on any given point $x$, merely by making a few random
% queries to $f$ (specifically, by querying $f(u)$ and $f(x + u)$ for a
% random choice of $u$).

% The BLR test was an early example of what is called  a ``property
% test'': a procedure for detecting the presence of global structure in
% an object, by only making a few local queries to it. Property tests
% were soon \anote{check!} developed for other classes of functions. 

% BLR phrased its result in terms of query access to an oracle. But
% there is an alternate formulation in terms of a two-player game. This
% is because of an idea known as ``oracularization'', which allows
% one to map from a strategy for a two-player game to an oracle and vice
% versa. This transformation meant that property tests could play a role
% in multiprover interactive proof systems, and indeed, the breakthrough
% result of BFL showing MIP = NEXP is at its heart a property testing
% result for \emph{multilinear} functions. 

% In generalizing these ideas to the quantum case, the natural starting
% point has been the view of self-tests as two-player games. A quantum
% self-test is typically understood as a game where success certifies
% that the entangled strategy used by the players has certain
% properties. Viewing a quantum self-test as a property test, the
% ``global object'' being tested is the entangled strategy, consisting
% of a shared state $\ket{\psi}$, and a collection of measurement
% operators indexed by the possible questions the players can receive in
% the game. The verifier performs ``local'' probes of this object in the
% sense that it always picks just one measurement for each player to
% perform, and observes a single outcome from each player. 


% Original motivation for property testing
% BLR
% 2-prover view
% Generalize to quantum: characterizing entangled strategies for
% nonlocal games


\section{Technical overview}\label{sec:technical-overview}


At a high level, our proof follows the approach of~\cite{BFL91}, and
it is useful to start by summarizing their analysis, which applies to
classical, \emph{deterministic} strategies. In this version of the test, the provers' strategy is described
by a ``points function,'' assigning a value in $\F_q$ to each point in
$\F_q^m$, and a
``lines function,'' assigning a low-degree polynomial to each
line in $\F_q^{m}$ queried in the test. From the assumption that the
points and lines functions agree at a randomly chosen point with high
probability, we would like to construct a global low-degree polynomial that has high agreement with
the points function. This is done by inductively constructing
``subspace functions'' defined on affine axis-aligned subspaces of increasing
dimension $k$. The base case is $k=1$, and is supplied by the lines
function. At each step, to construct the subspace function for a
subspace $S$ of dimension $k+1$, we pick $d+1$ parallel subspaces of
dimension $k$ that lie within $S$, and compute the unique degree-$d$
polynomial that interpolates them. The analysis shows that at each
stage, the function constructed through interpolation has high
agreement on average with the points function and lines function. In
the end, when we reach $k=m$ the ambient dimension, the subspace
function we construct is the desired global function.

For the sake of simplicity, it suffices to consider the case $d=1$,
i.e. multilinear functions. In this case, whenever we perform
interpolation, we need to combine $2$ parallel subspaces.

\paragraph{The classical zero-error case}

To start building intuition, it is useful to think about how to carry
out the above program in a highly simplified setting: the classical
zero-error case, for $m=3$. In this case, we assume that we have access to a
points function $f$ and lines function $g$ that \emph{perfectly} pass the BFL
test. Moreover, we will focus on the \emph{final} step of the
induction: thus, we assume that we have already constructed a set of
planes functions defined for every axis-parallel plane, that are
perfectly consistent with the line and point functions. In the final
step of the induction, our goal is to combine these planes functions
to create a single global function $h$
over all of $\F_q^{3}$ that is consistent with the points and lines functions. 

To do this, we will interpolate the planes as follows. Let us label
the 3 coordinates in the space $x, y, $and $z$, and consider planes
parallel to the $(x,y)$-plane. Each plane $S_z$ is specified by a
value of the $z$ coordinate:
\[ S_z = \{(x,y,z): (x, y) \in \F_q^2\}. \]
By the induction hypothesis, for every such plane $S_z$ there exists a
bilinear function $g_{z}: \F_q^{2} \to \F_q$ that agrees with the
points function. To construct a global multilinear function $h:
\F_q^3 \to \F_q$, we pick two distinct values $z_1 \neq z_2$, and ``paste''
the two plane functions $g_{z_1} $ and
$g_{z_2}$ together using \emph{polynomial
  interpolation}. Specifically, we define $h$ to be the unique
multilinear polynomial that interpolates between $g_{z_1}$ on the
plane $S_{z_1}$ and $g_{z_2}$ on the plane $S_{z_2}$.
\[ h(x,y,z) = \mathrm{interpolate}_{z_1, z_2}(g_{z_1}, g_{z_2}) = \left(\frac{z - z_2}{z_1 - z_2}\right) g_{z_1}(x,y) +
  \left(\frac{z - z_1}{z_2 - z_1}\right) g_{z_2}(x,y). \]

This procedure defines a global function $h$, and by the assumption
that $g_{z_1}$ and $g_{z_2}$ are multilinear, it follows that $h$ is
also a multilinear function. But why is $h$ consistent with the points
function? To show this, we need to consider the lines function on
lines parallel to the $z$ axis. Given a point $(x,y,z)$, let $\ell$ be
the line parallel to the $z$ axis through this point, and let $g_\ell$
be the associated lines function. By construction, $h$ agrees with
$g_\ell$ at the two points $z_1$ and $z_2$. But $g_\ell$ and the restriction
$h_{|\ell}$ of $h$ to $\ell$ are both linear functions, and hence if
they agree at two points, they must agree \emph{everywhere}. (This is
a special case of the \emph{Schwartz-Zippel lemma}, which says that if
two degree-$d$ polynomials agree at $d+1$ points, they must be equal.)
Thus, $h$
agrees with $g_\ell$ at the original point $z$ as well. By success in
the test, $g_\ell$ in turn agrees with the points function $f$ at
$(x,y,z)$, and thus, $h(x,y,z) = f(x,y,z)$. Thus, we have shown that
the global function $h$ is both multilinear and agrees with the points
function $f$ exactly.


\paragraph{Dealing with errors}
To extend the sketch above to the general case, with nonzero error,
requries some modifications. At the most basic level, we may consider
what happens when we allow for deterministic classical strategies that
succeed with probability less than $1$ in the test. Such strategies
may have ``mislabeling'' error: the points function
$f$ may be imagined to be a multilinear function that has been
corrupted at a small fraction of the points. This type of error is
handled by the analysis in~\cite{BFL91}. The main modification to the zero-error
sketch above is a careful analysis of the probability that the pasting
step produces a ``good'' interpolated polynomial for a randomly chosen
pair of planes $S_{z_1}, S_{z_2}$. This analysis makes use of the
Schwartz-Zippel lemma together with combinatorial properties of the
point-line test itself (e.g. the expansion properties of the question
graph associated with the test).

At the next level of generality, we could consider classical
\emph{randomized} strategies. Suppose we are given a randomized
strategy that succeeds in the test with probability $1 - \eps$. Any
randomized strategy can be modeled by first sampling a random seed,
and then playing a deterministic strategy conditioned on the value of
the seed. A success probability of $1 - \eps$ could have two
qualitatively different underlying causes: (1) on $O(\eps)$ fraction
of the seeds, the strategy uses a function which is totally corrupted,
and (2) on a large fraction of the seeds, the strategy uses functions
which are only $\eps$-corrupted. An analysis of randomized
strategies could naturally proceed in a ``seed-by-seed'' fashion, applying the
deterministic analysis of~\cite{BFL91} to the large fraction of ``good'' seeds (which
are each only $\eps$-corrupted), while giving up entirely on the
``bad'' seeds. 

In this document, we consider quantum strategies, which
have much richer possibilities for error. Nevertheless, we are able to
preserve some intuition from the randomized case, by working with
\emph{sub-measurements}: quantum measurements that do not always yield
an outcome. Working with a sub-measurement allows us to distinguish two
kinds of error: \emph{consistency error} (the probability that a
sub-measurement returns a wrong outcome) and \emph{completeness error}
(the probability that the sub-measurement fails to return an outcome at
all). Roughly speaking, the completeness error corresponds to the
probability of obtaining a ``bad'' seed in the randomized case, while
the consistency error corresponds to how well the strategies do on
``good'' seeds. 

The technique of using sub-measurements and managing the two types of
error separately goes back to~\cite{IV12}. That work developed a
crucial tool to convert between these two types of error called the
\emph{self-improvement lemma}, and in our analysis we make extensive
use of a refined version of this lemma (\Cref{thm:self-improvement-in-induction-section}), building on~\cite{Vid16,NV18b}. Essentially, the lemma says the
following: suppose that (a) the provers pass the test with probability
$1 - \eps$, and (b) there is a complete measurement $G_{g}$ whose
outcomes are low-degree polynomials that has
consistency error $\nu$: that is, $G$ always returns an outcome, but
has probability $\nu$ of producing an outcome $g$ that disagrees with the points measurement $A^u_a$ at a
random point $u$. Then there exists an
``improved'' sub-measurement  $H_{h}$ with consistency error $\zeta$
depending \emph{only} on $\eps$, and with completeness error
(i.e. probability of not producing an outcome at all) of $\nu
+\zeta$. Essentially, this lemma says we can always ``reset'' the
consistency error of any measurement we construct at intermediate
points in the analysis to a universal
function $\zeta$ depending only on the provers' success in the test,
at the cost of introducing some amount completeness error. Intuitively, one may
think of the action of the lemma as correcting $G$ on the portions of
Hilbert space where it is only mildly corrupted, while ``cutting out'' the
portions of Hilbert space where $G$ is too corrupted to be
correctable. In some sense, this lemma is the quantum analog of the
idea of identifying ``good'' and ``bad'' random seeds in the classical
randomized case. The proof of the lemma uses a semidefinite program
together with the combinatorial facts used in~\cite{BFL91}
(specifically, expansion of the question graph of the test, and
the Schwartz-Zippel lemma).

Armed with the self-improvement lemma, we set up the following
quantum version of the BFL
induction loop: for $k$ running from $1$ to $m$, we construct a family
of subspace measurements $G^{S}_{g}$ that returns a low-degree
polynomial $g$ for every axis-aligned affine subspace $S$ of dimension $k$.
\begin{enumerate}
  \item By the induction hypothesis, we know that there exists a
    measurement  $G^{S}_{g}$ for every $k$-dimensional subspace, which
    has consistency error $\delta(k)$ with the points
    measurement. (For the base case $k = 1$, this is the lines
    measurement from the provers' strategy). 
  \item We apply the self-improvement lemma to these measurements, yielding
    sub-measurements $\hat{G}^{S}_{g}$ that have consistency error
    $\zeta$ independent of $\delta(k)$, and completeness error $\kappa(k) = \delta(k) + \zeta$.
  \item For each subspace $S$ of dimension $k+1$, we construct a
    pasted sub-measurement, by performing a quantum version of the
    classical interpolation argument: we define the pasted
    sub-measurement by sequentially measuring several subspace
    measurements corresponding to parallel $k$-dimensional subspaces,
    and interpolate the resulting outcomes. This pasted sub- measurement has consistency slightly worse than $\zeta$, and completeness error which
    is slightly worse than $\kappa(k)$. It is at this step that it is crucial to treat the two types
    of error separately: in particular, we need the consistency error
    to be low to ensure that the interpolation produces a good result.
  \item We convert the resulting sub-measurement into a full
    measurement, by assigning a random outcome whenever the
    sub-measurement fails to yield an outcome. This measurement will
    have consistency error $\delta(k+1)$ which is larger than $\delta(k)$
    by some additive factor.
  \end{enumerate}
  At the end of the loop, when $k = m$, we obtain a single measurement that returns
  a global polynomial as desired.
\paragraph{The diagonal lines test}
An important element of the test we analyze in this document, which is
unnecessary in the classical case, is the diagonal lines test. The purpose
of this test is to certify that the points measurements used by the provers
approximately commute on average over all pairs of points $(x, y)$ in
$\F_q^m$---something which is automatically true in the classical case. This is done by asking one prover for a polynomial defined
on the line going through $x$ and $y$, and the other for the function
evaluation at either $x$ or $y$. The line through $x$ and
$y$ will not in general be axis-parallel; we refer to these general
lines as ``diagonal.''

The commutation guarantee plays an important role in our analysis of
the test, and it is an interesting question whether it is truly
necessary to test it directly with the diagonal lines test; might it
not automatically follow from success in the axis-parallel lines test? An interesting contrast can be drawn to
the Magic Square
game~\cite{mermin1990simple,peres1990incompatible,aravind2002simple},
in which questions are either cells or axis-parallel lines in a $3 \times 3$ square grid.
This has the same question distribution as
the axis-parallel line-point test over $\F_3^2$ (although the answers in the Magic Square game are strings in~$\F_2$ rather than~$\F_3$).
For the Magic Square game,
``points'' measurements along the same axis-parallel ``line'' commute,
but points that are not axis-aligned do \emph{not} commute: indeed,
for the perfect strategy, they anticommute. Despite this example,
we know that commutation between all pairs of points \emph{can} be deduced from the
axis-parallel lines test alone, rending the diagonal lines test unnecessary,
at least in the case of the bivariate ($m = 2$) multilinearity test (the low individual degree test when $d = 1$).
On the other hand, we know of a quantum strategy using noncommuting measurements
which succeeds with probability~$1$ in the $m = 2$, $d = 2$, $q = 4$ low individual degree test.
Whether this counterexample can be extended to larger~$q$ remains an open question.

\paragraph{Organization}
The rest of this document is organized as follows. In
\Cref{sec:prelims} we review some preliminaries concerning finite
fields, polynomials, and quantum measurements. In
\Cref{sec:making-measurements-projective}, we present two tools for
making quantum measurements projective, which are used in our
analysis. In \Cref{sec:induction}, we give the inductive argument
proving the main theorem. In \Cref{sec:expansion,sec:variance} we show
some properties of the hypercube graph and families of measurements
indexed by points on the hypercube, which are used in
\Cref{sec:self-improvement} to prove the self-improvement lemma. In
\Cref{sec:commutativity-points,sec:g-comm}, we show commutativity
properties of the measurements constructed in the induction, and
finally in \Cref{sec:ld-pasting} we analyze the pasting step of the induction.


\paragraph{Acknowledgments}

We thank Lewis Bowen for pointing out typos and a few minor errors in a previous version.
We also thank Madhu Sudan for help with references on the classical low individual degree test.

% \ignore{
% \section{The test}

% \begin{definition}
%   A strategy $(\psi, A, B, L)$ is $(\eps, \delta, \gamma)$-good if it
%   succeeds in the consistency test with probability at least $1 -
%   \delta$, in the axis-parallel line test with probability at least $1
%   - \eps$, and in the diagonal line test with probability $1 - \gamma$.
% \end{definition}
% \begin{theorem}\label{thm:main}
% Suppose $(\psi, A, B, L)$ forms an $(\eps, \delta,\gamma)$-good
% strategy for the multilinearity test.  Then there exists a state
% $\ket{\rmaux}$ and a measurement $G \in \multimeas{n}{q}$ such that on
% the state $\ket{\tilde{\psi}} = \ket{\psi} \ot
% \ket{\rmaux}_{\reg{aux_A}} \ot \ket{\rmaux}_{\reg{aux_B}}$, the
% following consistency relation holds:
% \begin{equation}\label{eq:example}
% (A^w_a \ot I_{\reg{aux_A}}) \otimes I_{\reg{Bob}} \simeq_{\xi} I_{\reg{Alice}} \otimes G_{[g(w) = a]},
% \end{equation}
% where $\xi = XXX$.
% \end{theorem}

% To prove the main theorem, we will construct measurements defined over
% larger and larger axis-aligned subspaces of $\F_q^m$. The following
% lemma states the properties of the intermediate measurements we construct.
% \begin{lemma}
% Suppose $(\psi, A, B, L)$ forms an $(\eps, \delta,\gamma)$-good
% strategy for the multilinearity test, and let $1 \leq k \leq
% m$. \ainnote{Add in some text about the universal aux state for Naimark.}  Then
% for every $u_{<k} \in \F_q^{k}$ there exists a projective
% submeasurement $G^{u_{<k}} \in
% \multimeas{m - k+1}{q}$ and convex\anote{check!} functions $\kappa(k,
% u_{<k})$ and $\eta(k, u_{<k})$ such that
% \begin{enumerate}
% \item (Completeness):\label{item:main-induction-completeness} If $G^{u_{<k}} =  \sum_g G^{u_{<k}}_g$, then
%   \[\bra{\psi} G^{u_{<k}} \otimes I  \ket{\psi} \geq 1 - \kappa(k, u_{<k}). \]
% \item (Consistency with~$A$):\label{item:main-induction-A-consistency} On average over points $w$ whose first $k-1$ coordinates are fixed to
% $u_{<k}$,
% \begin{equation}
% A^w_a \otimes I \simeq_{\eta(k, u_{<k})} I \otimes G^{u_{<k}}_{[g(w) = a]}.
% \end{equation}
% \item (Self-consistency):\label{item:main-induction-self-consistency}
%   \[ G^{u_{<k}}_{g} \ot I \approx_{\eta(k, u_{<k})} I \ot G^{u_{<k}}_{g}. \]
% \item (Boundedness):\label{item:main-induction-boundedness} For each $u_{<k}$, there exists a positive
%   semidefinite matrix $Z^{u_{<k}}$ such that
%   \[ \bra{\psi} (I - G^{\bu_{<k}}) \ot Z^{\bu_{<k}}
%     \ket{\psi} \leq \eta(k, u_{<k}), \]
%   and for each $h \in \multi{m-k+1}{q}$,
%   \[ Z^{u_{<k}} \geq \left( \E_{\bv \in \F_{q}^{m-k+1}} A^{(u_{<k},
%         \bv)}_{h(\bv)} \right). \]
% \item (Average errors):\label{item:main-induction-avg-error}
%   \begin{align*}
%     \E_{\bu_{<k}} \kappa(k, \bu_{<k}) &\leq 300 (m-k+1) \cdot \left( \left(
%                                         \frac{dm}{q} \right)^{1/8} +
%                                         (m^2 \eps)^{1/32} + (m
%                                         \delta)^{1/32} + (m^2
%                                         \gamma)^{1/32}  \right. \\
%     &\qquad\qquad + \left. 2
%                                         \cdot \left(m \cdot \left(m \eps +
%                                         m \delta +
%                                         \frac{1}{q}\right)\right)^{1/64}\right)\\
%     \E_{\bu_{<k}}
%     \eta(k, \bu_{<k}) &\leq 100 \cdot \left(m \cdot \left(m  \eps
%                         + m \delta + \frac{1}{q} \right) \right)^{1/2}.
%   \end{align*}
% \end{enumerate}
% \label{lem:main-inductive}
% \end{lemma}
% \begin{proof}[Proof of \Cref{lem:main-inductive}]
%   We prove the lemma by downwards induction on $k$. The base case is
%   $k = m$, in which case the prefix $u_{<m}$ fixes the first $m-1$
%   coordinates. In this case, we set
%   \[ G^{u_{<m}}_{g} = B^{u_{<m}}_{g}. \]
%   Since $B$ is a complete projective measurement, $G^{u_{<m}}$ is one
%   as well, and we may trivially satisfy the boundedness property
%   (\Cref{item:main-induction-boundedness}) with
%   with sufficiently large $Z$ matrix (for instance, setting $Z^{u_{<m}} = I$ does
%   the job). Since the given strategy succeeds in the axis-parallel lines test
%   with probability $1 - \delta$ and the self-consistency test with
%   probability $1 - \eps$, it follows that
%   \Cref{item:main-induction-A-consistency,item:main-induction-self-consistency}
%   hold with the following error bounds:
%   \begin{align*}
%     A^{w}_{a} \otimes I &\simeq_{m \cdot \delta} I \otimes
%                           G^{u_{<m}}_{[g(w) = a]} \\
%     G^{u_{<m}}_{g} \ot I &\simeq_{m \cdot \eps} I \otimes G^{u_{<m}}_{g}.
%   \end{align*}
%   It is clear from these expressions that
%   \Cref{item:main-induction-avg-error} is also satisfied.
  
%   We will now show the inductive step. Assume the lemma holds for
%   $k+1$. This means that for each prefix $u_{<k} \in \F_q^{k - 1}$,
%   there exists a family of measurements $G^{u_{<k}, x}_{g}$ indexed
%   by elements $x \in \F_q$  satisfying the conclusions of the
%   lemma. Define
%   \begin{align*}
%     \kappa'(k, u_{<k}) &= \E_{\bx} \kappa(k+1, (u_{<k}, x))\\
%     \eta'(k, u_{<k}) &= \E_{\bx} \eta(k+1, (u_{<k}, x))\\
%   \end{align*}
%   We will now apply \Cref{thm:pasting} separately to the
%   $G^{u_{<k}, x}_{g}$ measurements for each choice of $u_{<k}$.

%   Note
%   that for every fixing of $u_{<k}$, the strategy $(\psi, A, B, L)$
%   defines a strategy for the $u_{<k}$-restricted\anote{define!} test
%   which is $(\eps_{u_{<k}}, \delta_{u_{<k}}, \gamma_{u_{<k}})$-good, where
%   \begin{align}
%     &\E_{\bu_{<k}} \eps_{u_{<k}} \leq m \cdot \eps, \quad
%     \E_{\bu_{<k}} \delta_{u_{<k}} \leq m \cdot \delta, \quad
%     \E_{\bu_{<k}} \gamma_{u_{<k}} \leq m \cdot \gamma. \label{eq:subspace-avg-error}
%   \end{align}

%   From \Cref{thm:pasting}, we deduce the
%   existence of measurements $H^{u_{<k}}_{h} \in \multisub{m-k+1}{q}$ with the following
%   properties:
%   \begin{enumerate}
%   \item (Consistency with~$A$): On average over points $w \in \F_q^m$
%     with prefix $u_{<k}$,
%     \[ A^{w}_{a} \ot I \simeq_{\nu(k, u_{<k})} I \ot H^{u_{<k}}_{[h(w)
%         = a]}. \]
%   \item (Completeness)
%     \[ \sum_{h} \bra{\psi} H^{u_{<k}}_{h} \ot I \ket{\psi} \geq 1 -
%       \kappa'(k, u_{<k}) - \nu(k, u_{<k}). \]
%   \end{enumerate}
%   where
%   \begin{align*}
%     \nu(k, u_{<k}) &= \nu(d, m - k + 1, q, \eps_{u_{<k}},
%                      \delta_{u_{<k}}, \gamma_{u_{<k}}, \eta'(k, u_{<k})) \\
%                    &= 100 \cdot \left( \left(\frac{d
%                      (m-k+1)}{q}\right)^{1/8} + (m \cdot
%                      \eps_{u_{<k}})^{1/32} + \delta_{u_{<k}}^{1/32} +
%                      (\gamma_{u_{<k}} \cdot (m - k))^{1/32} + \eta'(k, u_{<k})^{1/32}\right).
%   \end{align*}

%   To construct the measurements $G^{u_{<k}}$ and PSD matrices
%   $Z^{u_{<k}}$ in the conclusion of the
%   lemma, we now apply \Cref{thm:self-improvement} to $H^{u_{<k}}$ to
%   obtain an ``improved'' measurement $\tilde{H}^{u_{<k}}$ and a
%   collection of PSD matrices $\tilde{Z}^{u_{<k}}$. We set
%   $G^{u_{<k}}$ to be the Naimark dilation (\Cref{thm:naimark}) of
%   $\tilde{H}^{u_{<k}}$, and $Z^{u_{<k}} = \tilde{Z}^{u_{<k}} \ot I_{\reg{aux}}$. Note that $G^{u_{<k}}$ is a projective
%   submeasurement by construction, and moreover that by
%   \Cref{cor:strat-naimark}, $G$ and $Z$ satisfy all of the conclusions of
%   \Cref{thm:self-improvement} on the dialted state with the same error bounds as
%   $\tilde{H}^{u_{<k}}$\anote{explain this better, in particular the
%     boundedness condition}.  We thus
%   obtain that
%   \Cref{item:main-induction-completeness,item:main-induction-A-consistency,item:main-induction-self-consistency,item:main-induction-boundedness}
%   of the lemma
%   hold with
%   \begin{align*}
%     \kappa(k, u_{<k}) &= \kappa'(k, u_{<k}) + 2\cdot \nu(k, u_{<k}) +
%                         \zeta(d, m-k+1,q, \eps_{u_{<k}},
%                         \delta_{u_{<k}}, \gamma_{u_{<k}}) \\
%     \eta(k, u_{<k}) &= \zeta(d, m-k+1,q, \eps_{u_{<k}},
%                         \delta_{u_{<k}}, \gamma_{u_{<k}}),
%   \end{align*}
%   where
%   \[ \zeta(d,m,q,\eps, \delta, \gamma) = 100 \cdot \left(m \cdot
%       \left(\eps + \delta + \frac{1}{q} \right) \right)^{1/2}. \]
%   It remains to show \Cref{item:main-induction-avg-error} holds for
%   the errors computed above. This will requiring averaging the error
%   over the choice of $u_{<k}$ and using the convexity properties of
%   the relevant functions. First, we show the desired bound for the
%   average of $\eta$:
%   \begin{align*}
%     &\E_{\bu_{<k}} \eta(k, \bu_{<k}) \\
%     =~&\E_{\bu_{<k}} \zeta(d, m-k+1,q, \eps_{u_{<k}},
%         \delta_{u_{<k}}, \gamma_{u_{<k}}) \\
%     =~&\E_{\bu_{<k}} \left( 100 \cdot  \left((m-k+1) \cdot
%       \left(\eps_{\bu_{<k}} + \delta_{\bu_{<k}} + \frac{1}{q} \right)
%         \right)^{1/2} \right) \\
%     \leq~&100 \cdot \left((m-k+1) \cdot \left( \E_{\bu_{<k}}
%            (\eps_{\bu_{<k}} + \delta_{\bu_{<k}}) + \frac{1}{q} \right)
%            \right)^{1/2} \tag{by convexity} \\
%     \leq~&100 \cdot \left(m \cdot \left(m \cdot \eps + m \cdot \delta
%            + \frac{1}{q} \right)\right)^{1/2}. \tag{by \Cref{eq:subspace-avg-error}}
%   \end{align*}
%   Now we show the desired bound for the average of $\kappa$:
%   \ainnote{TODO: compute the average for $\kappa$.}
%   \begin{align*}
%     &\E_{\bu_{<k}} \kappa(k, \bu_{<k}) \\
%     =~&\E_{\bu_{<k}} \left( \kappa'(k, \bu_{<k}) + 2 \cdot \nu(k,
%         \bu_{<k}) + \zeta (d, m - k+1, q, \eps_{\bu_{<k}},
%         \delta_{\bu_{<k}}, \gamma_{\bu_{<k}})\right) \\
%     =~&\E_{\bu_{<(k+1)}} (\kappa(k+1, \bu_{<(k+1)} ) + \E_{\bu_{<k}}
%         \left( 2 \cdot \nu(k,
%         \bu_{<k}) + \zeta (d, m - k+1, q, \eps_{\bu_{<k}},
%         \delta_{\bu_{<k}}, \gamma_{\bu_{<k}})\right) \\
%     \leq~&\E_{\bu_{<(k+1)}} (\kappa(k+1, \bu_{<(k+1)} ) +
%            \E_{\bu_{<k}} 3 \cdot \nu(k, \bu_{<k})
%   \end{align*}
% \end{proof}

% \begin{proof}[Proof of \Cref{thm:main}]
% Apply \Cref{lem:main-inductive}
% setting $k$ to be equal to $1$ to obtain a projective submeasurement
% $G \in \multisub{m}{q}$, and complete $G$ arbitrarily. Then the
% conclusion of \Cref{thm:main} holds with
% \begin{align*}
%   \xi &= \kappa(k=1, u_{<k} =
%         \emptyset) + \eta(k=1, u_{<k} = \emptyset)\\
%   &\leq  300 \cdot m\cdot \left( \left(
%                                         \frac{dm}{q} \right)^{1/8} +
%                                         (m^2 \eps)^{1/32} + (m
%                                         \delta)^{1/32} + (m^2
%                                         \gamma)^{1/32}  +  2
%                                         \cdot \left(m \cdot \left(m \eps +
%                                         m \delta +
%                                         \frac{1}{q}\right)\right)^{1/64}\right)\\
%       &\qquad + 100 \cdot \left(m \cdot \left(m  \eps
%                         + m \delta + \frac{1}{q} \right) \right)^{1/2}.
% \end{align*}
% \end{proof}
% }
%%% Local Variables:
%%% mode: latex
%%% TeX-master: "multilinearity"
%%% End:
